\minitopic{从直觉到论证：四种推理任务}初学者常把所有理由都叫作“证明”，但认识论中的支持关系至少有四类。第一，\textbf{演绎推理}要求前提为真时结论不可能为假。它的强度最高，却不能凭空增加经验内容。第二，\textbf{归纳推理}从样本、频率或过去规律支持未观察结论，结论可能为假。第三，\textbf{溯因推理}或最佳解释推理，从现象出发比较能解释它们的假设。第四，\textbf{证言接受}依赖说话者、传播链和听者的认识地位，不能简单归入个人观察。 区分推理类型能防止错置标准。要求医学预测具有演绎确定性，会把一切可错科学排除；仅因某解释“说得通”就接受它，又会忽略竞争解释和先验可信度。评价时应先问：这一推理声称提供必然保证、概率提高、解释优势，还是可被击败的信任权利？只有弄清承诺，反例才击中目标。 演绎有效性也不保证前提合理。论证“所有梦中经验都不可靠；当前经验可能是梦；所以当前经验不可靠”即使形式上可整理为有效，第一前提仍需辩护。相反，一个结论很可能为真，也不能挽救无效推理。认识论评价同时关心形式联系、前提地位和主体是否适当基于这些前提。

\minitopic{论证图的制作}阅读时可给每个陈述编号，并用箭头标记支持。两个前提共同支持结论时，要用同一括号连接；各自独立支持时，则画两条箭头。反驳也有不同对象：攻击某个前提，攻击从前提到结论的联系，提出相反证据，或说明即使结论成立也不能回答原问题。 以“知觉有时出错，所以知觉从不产生知识”为例：前提只是偶称命题，结论是全称命题，中间需要“只要一种来源可能出错，它就从不产生知识”。把桥梁前提写出后，讨论从笼统的赞成或反对，转向可错论能否接受的精确争点。 论证图还会暴露循环。如果“这份报告可信，因为它来自可靠机构”，而“机构可靠，因为它总发布可信报告”，两条理由彼此支撑却没有独立输入。这不表示任何相互支持都是恶性循环；长期记录、外部审计和可重复预测可以给循环系统增加世界约束。

\minitopic{概念工程：我们只是在描述用法吗}概念分析可有三种目标。描述性分析试图说明日常“知道”如何使用；解释性分析寻找在理论中统一现象的结构；改良性或\term{概念工程}{conceptual engineering}则问，为了研究或公共实践，我们应当使用怎样的概念。三者的证据不同。日常语言调查对描述性目标重要，却不能单独决定哪种概念最公正或最有解释力。 例如，日常群体知识归属可能摇摆：“委员会知道风险”有时只是说报告已进入档案，有时表示组织能据此行动。描述者记录这种用法；解释者区分信息存储、正式接受与行动能力；改良者可能建议，在安全监管中只有能够触发责任流程的信息才称为组织知识。建议不是发现隐藏词义，而是基于实践目标重设规范。 概念工程也有风险。若哲学家把“知识”重新规定成自己理论最容易满足的状态，就可能逃避原问题；若改革完全脱离现有用法，又难以指导真实交流。因此合理改良需要说明现有概念造成何种混乱，新概念改善什么，以及转换成本和可能排除谁。

\minitopic{自然主义与规范性}认识论是否应成为心理学的一部分？自然主义者强调，人类实际如何形成信念、感知系统在什么环境可靠、偏差怎样产生，不能只靠先验思辨。蒯因提出把传统的基础奠定计划转化为研究感觉输入如何导致理论输出的经验项目\parencite{quine1969}。这一转向推动认知科学与认识论对话。 反对意见是，心理学描述人事实上怎样相信，不直接回答应当怎样相信。若实验发现人普遍受确认偏误影响，结论通常是这种倾向需要校正，而不是因为普遍所以合理。这是“是”与“应当”的间隙。自然主义可以回应：规范目标仍是获得真理或成功行动，经验科学帮助我们找出实现目标的可靠策略；规范性没有被消除，而是与可行机制连接。 另一个问题是可靠性环境相对。依赖面孔识别在熟悉群体中可能可靠，在跨群体环境中显著下降；启发式在信息稀缺时有效，在广告操纵下却系统失灵。经验研究不仅给过程一个总准确率，还应揭示能力—环境配合。这样的自然化结果反过来修改哲学案例中“正常条件”的含义。

\minitopic{实验哲学能告诉我们什么}实验哲学收集不同人群对思想实验的判断，研究措辞、顺序、文化和个体差异。它可以检验哲学家声称“人人都会判断”的经验前提，也能发现一个案例实际触发的是道德责备、概率估计还是知识概念。调查不是用多数票决定真理，而是研究概念判断这类证据的分布与机制。 设计调查至少要控制四点。第一，参与者是否理解技术词；第二，翻译是否保持情境而非逐词对应；第三，问题是否把“知道”“确定”“有理由”等评价混在一起；第四，案例顺序是否产生对比效应。若不同群体判断不同，还需说明差异稳定、可复制，并与哪些背景变量相关。 实验结果对不同理论目标的意义不同。若目标是描述普通语言，广泛而稳定的差异非常重要；若目标是提出规范性知识概念，多数直觉可能本身受偏见影响；若思想实验只用于展示理论后果，即使直觉不一致，它仍能说明条件组合。把“发现直觉差异”直接等同于“反驳扶手椅哲学”过于迅速。

\minitopic{认识论与伦理学的边界}相信是否受意志控制？我们似乎不能仅因有利而立即相信明天晴天，却能选择查证、注意哪类信息、接近哪些来源。认识责任更多落在间接控制：是否培养能力、是否忽视反证、是否让愿望支配搜索。因而“你不该相信$p$”可能批评信念本身与证据不相称，也可能批评导致信念的可控行为。 认识规范与道德规范可能冲突。医生保护隐私可能限制公开证据；记者为避免伤害会延迟发布已核实信息；朋友有时认为信任本身是关系义务。认识论不能假定追求每个真相永远优先，但也不能因道德目的良好就把愿望当证据。正确做法是把两类理由分别列出，再讨论行动选择，而不是改变命题的证据支持。

\minitopic{综合案例：校园匿名指控}校园论坛出现匿名帖，指控一名教师篡改成绩。帖子包含一张局部截图，迅速被多个账号转发。我们可以逐层分析。真值问题是成绩是否被篡改；证据问题是截图是否完整、账号是否独立、教务记录怎样；证言问题是匿名是否有保护理由、发布者是否有直接接触；社会问题是权力差异会不会使实名举报代价过高；责任问题是转发者应采取何种核验。 两种极端都不合理：匿名所以必假，忽略了弱势者可能只有匿名渠道；指控严重所以应立即相信，又把严重性误当作真实性证据。较合适的态度可能是：认为指控值得调查，对事实暂缓判断，保存材料并交由有程序保障的独立机构核验。这里“认真对待”不等于“已经知道为真”。 若后续证实成绩确被改动，早期轻率转发者不会因此自动获得知识；真值不能回溯修复当时根据。若调查证明截图被误读，早期谨慎要求调查者也不因此不理性。这个案例综合展示真值、确证、行动和责任评价为何必须分离。

\begin{tcolorbox}[breakable,colback=white,colframe=blue!30!black,title={方法清单},fonttitle=\bfseries]
面对新问题，依次写下：目标命题；评价类型；信息来源；推理类型；未明说的桥梁前提；可能的击败者；最小反例；竞争理论的代价；需要经验调查还是概念分析；在不确定性下可以采取的行动。
\end{tcolorbox}
